\documentclass[runningheads]{llncs}
\usepackage[T1]{fontenc}
\usepackage{graphicx}
\usepackage{booktabs}
\usepackage{amsmath}
\usepackage{amssymb}
\usepackage{url}
\usepackage{xcolor}
\usepackage{hyperref}
\begin{document}
\title{Design and Evaluation of an Explainable Hybrid Intrusion Detection and Risk Prioritization Framework for IoT Networks}
\titlerunning{Explainable Hybrid IDS for IoT Networks}
\author{[Author Name]}
\authorrunning{[Author Name]}
\institute{[Department and Institution]\\\email{[email address]}}
\maketitle

\begin{abstract}
The growing deployment of Internet of Things (IoT) devices has expanded the network attack surface while introducing constraints that complicate security monitoring, including heterogeneous protocols, limited device resources, and large volumes of traffic. This paper presents an offline, flow-based hybrid intrusion detection and risk prioritization framework that combines a supervised Random Forest (RF) classifier with a benign-trained Isolation Forest (IF) anomaly detector. The two detectors are joined using an OR decision rule: a flow is flagged when either component indicates malicious activity. The framework is evaluated on a memory-conscious, category-wise sample of CICIoT2023 using source-CSV group separation to reduce the risk of train--test contamination. The sampled dataset contains 487,091 records, 39 numeric features, and 309 source-file groups. On the held-out test partition, RF achieves 88.63\% accuracy, 85.76\% attack precision, 86.68\% attack recall, and an attack F1-score of 86.22\%. IF achieves 96.09\% attack precision but only 44.59\% recall. The hybrid rule raises recall marginally to 86.81\%, while reducing precision to 85.26\% and increasing the false-positive rate from 10.02\% to 10.45\%. Permutation importance identifies HTTPS, Max, Number, Time\_To\_Live, and Header\_Length among the most influential features. A probability--anomaly composite score is also used to rank alerts for analyst review. The results show that anomaly detection can provide a complementary signal, but a simple OR fusion does not improve the overall F1-score in this experiment. The paper therefore emphasizes transparent trade-offs, group-aware evaluation, and operational limitations rather than claiming superiority over commercial security products.
\keywords{Internet of Things \and Intrusion detection \and Random Forest \and Isolation Forest \and Hybrid detection \and Risk prioritization \and CICIoT2023}
\end{abstract}

\section{Introduction}
The Internet of Things (IoT) connects sensors, cameras, gateways, appliances, and industrial controllers to networks and cloud services. These deployments support smart homes, manufacturing, healthcare, logistics, and infrastructure monitoring. Their diversity creates security challenges: devices may have limited processing capacity, weak credentials, exposed management interfaces, and infrequent firmware updates. Compromised devices can participate in DDoS activity, scan neighbouring systems, or provide footholds for lateral movement.

Firewalls and signature-based network intrusion detection systems (NIDS), including Snort and Suricata, remain essential for enforcing policies and detecting known patterns. However, signatures require maintenance and may not cover novel behaviours. Machine-learning (ML) detectors can learn statistical traffic patterns, but may generate false alarms, inherit dataset bias, and provide limited explanations.

This work evaluates a lightweight hybrid design combining supervised classification and benign-only anomaly detection. An OR rule flags a flow if either model predicts an attack. A separate score combines RF attack probability and normalized anomaly magnitude to rank alerts for review.

The study uses CICIoT2023, collected from an IoT testbed of 105 devices and 33 attacks [1]. Because the dataset exceeds 46 million records, chunk-based sampling is used. Source CSV files are separated into fit, calibration, and test groups to reduce within-file leakage, although this does not guarantee independence across captures or devices. The contributions are a memory-conscious pipeline, group-aware evaluation of RF/IF/OR fusion, analysis of detection trade-offs, and a transparent risk-ranking heuristic. This is a research prototype, not a production NIDS or a direct commercial-product benchmark.

\section{Literature Review}
\subsection{IoT intrusion detection and benchmark datasets}
Neto et al. introduced CICIoT2023, a benchmark built from 105 real IoT devices and 33 attacks spanning DDoS, DoS, reconnaissance, web attacks, brute force, spoofing, and Mirai [1]. It provides packet captures and flow-level CSV features; this study uses only the latter. Evaluation protocol matters: random row splitting can distribute related flows from one capture across training and test sets. Grouping by source CSV reduces this risk, but separate files may still share devices, campaigns, or collection conditions. It therefore strengthens within-dataset evaluation without proving generalization to arbitrary IoT networks.

\subsection{Signature-based and machine-learning approaches}
Snort and Suricata use rules and protocol inspection to identify known patterns, offering readable alerts and mature workflows. Their coverage depends on maintained signatures and suitable behavioural rules. Commercial NDR products may combine signatures, analytics, and threat intelligence, but their proprietary capabilities are not directly comparable with this offline benchmark; no superiority claim is made.

Random Forest (RF) combines decision trees [3] and models nonlinear relationships in tabular flow features without extensive scaling. It can provide global feature rankings, but relies on labelled examples and may degrade under distribution changes or learn dataset artefacts when evaluation is weak.

Isolation Forest (IF) isolates observations through random feature and split selection; anomalies tend to have shorter isolation paths [2]. Fitting it on benign traffic avoids requiring examples of every attack. However, rare legitimate behaviour may be flagged and subtle attacks may appear normal, making threshold selection central to the precision--recall trade-off.

\subsection{Hybrid detection, explainability, and alert prioritization}
Hybrid IDS designs combine supervised recognition of known patterns with anomaly detection for deviations. Logical OR fusion may increase recall but can also increase false positives; weighted fusion and context-aware thresholds are alternatives. The choice should be evaluated against operational costs, not accuracy alone.

Global permutation importance measures how a chosen score changes when a feature is shuffled [4]. It provides a model-agnostic feature ranking, but not a local explanation for each alert; correlated features can share importance. SHAP is not implemented here. Thus, explainability in this work means global feature transparency and visible score components.

Alert triage is also important when teams cannot investigate every alert immediately. Prioritization may use attack probability, anomaly magnitude, asset criticality, and context. This prototype uses only the first two and is not a calibrated estimate of loss or incident severity.

\section{Technology Used to Solve the Problem}
\subsection{Framework overview}
The proposed framework has five stages: (1) flow CSV ingestion and category-wise sampling; (2) numeric feature preparation; (3) group-based data partitioning; (4) RF and IF model training and hybrid decision; and (5) evaluation, explanation, and risk ranking. Fig.~\ref{fig:pipeline} summarizes the workflow.

\begin{figure}[t]
\centering
\fbox{\parbox{0.92\linewidth}{\centering
CICIoT2023 CSV files\\[2mm]
$\downarrow$ chunked ingestion and category-wise reservoir sampling\\[2mm]
Numeric flow features + source-file group identifier\\[2mm]
$\downarrow$ group-disjoint fit / calibration / test partitions\\[2mm]
\begin{tabular}{ccc}
Random Forest & Isolation Forest & Benign-only fit\\
Supervised & Anomaly detection & Calibration threshold
\end{tabular}\\[2mm]
$\downarrow$ RF prediction OR IF anomaly flag\\[2mm]
Metrics + permutation importance + risk-ranked alerts
}}
\caption{Architecture of the proposed offline hybrid intrusion detection workflow.}
\label{fig:pipeline}
\end{figure}

\subsection{Dataset and memory-conscious sampling}
CICIoT2023 contains flow-level statistical features extracted from traffic generated in an IoT testbed [1]. The full collection contains more than 46 million records. Loading all CSV data into memory was not feasible on the target 16-GB laptop. The implementation therefore reads each file in chunks of 100,000 rows and maintains a reservoir sample for each parent-folder category. The sampling caps are 200,000 benign records and up to 10,000 records per attack category. This produced 487,091 sampled records from 309 CSV files, represented by 39 numeric features. The sample includes 200,000 benign records and 287,091 attack records.

This is a category-capped sample, not a prevalence-preserving sample of the original network. The resulting attack proportion (58.94\%) is substantially different from the original dataset distribution. Thus, precision, accuracy, and false-positive rate measured here should not be interpreted as expected deployment performance under natural traffic prevalence. The sampling design improves coverage of attack categories and makes the experiment computationally manageable, but changes the class prior.

\subsection{Preprocessing and group-aware partitioning}
The parent directory name supplies the attack category. Records in the \texttt{Benign\_Final} directory are assigned label 0; all other categories are assigned label 1. The category and target are retained as metadata and excluded from model inputs. The source CSV path is retained as a group identifier and also excluded from features. Identifier-like columns (such as flow IDs, IP addresses, ports, and timestamps) are excluded. Remaining columns are converted to numeric values; infinite values are replaced with missing values, all-empty columns are dropped, and missing values are handled by a median imputer in the model pipeline.

A GroupShuffleSplit-based procedure assigns complete source CSV files to three disjoint partitions: model fitting, anomaly-threshold calibration, and final testing. The executed run placed 215 groups in fit, 47 in calibration, and 47 in test. These correspond to 69.6\%, 15.2\%, and 15.2\% of groups, respectively. Because CSV files differ in size, the record proportions were 52.3\% fit (254,594 records), 28.0\% calibration (136,431), and 19.7\% test (96,066). The group-overlap assertion passed. Reporting both group and row proportions is essential; the split should not be described as a 70/15/15 row split.

\subsection{Random Forest classifier}
The RF model is a supervised binary classifier trained on the fit partition. It outputs a class label and an estimated attack probability $p_{RF}(x)$. Median imputation is included in the preprocessing pipeline. The model uses 120 trees, square-root feature subsampling, a minimum leaf size of two, balanced-subsample class weighting, and random seed 42; two CPU workers are used. The RF model is the primary detector and supplies the probability component of the risk score.

\subsection{Isolation Forest anomaly detector}
The IF detector is fitted only on benign examples from the fit partition. Its anomaly decision threshold is calibrated using benign examples in the separate calibration partition. Let $a(x)$ denote the anomaly score and $\tau$ the calibrated threshold. The anomaly flag is
\begin{equation}
I_{IF}(x)=\begin{cases}1,&a(x)>\tau,\\0,&\text{otherwise.}\end{cases}
\end{equation}
The anomaly component is not trained to recognize attack labels. It instead models the distribution of benign feature vectors using 120 isolation trees and a fixed random seed. The calibration threshold is the 95th percentile of anomaly scores on benign calibration records, corresponding to a target calibration FPR of 5\%. This distinction explains why IF can have high precision but lower recall: many attacks may not be sufficiently anomalous under the learned benign profile.

\subsection{Hybrid decision and risk score}
The hybrid prediction uses a logical OR:
\begin{equation}
\hat{y}_{H}(x)=\mathbb{1}[\hat{y}_{RF}(x)=1\;\lor\;I_{IF}(x)=1].
\end{equation}
This rule prioritizes sensitivity: either detector can raise an alert. The expected trade-off is that recall may increase while false positives also increase.

For triage, the notebook computes a normalized anomaly value $a_n(x)$ over the test scores and a composite score:
\begin{equation}
R(x)=100\left(0.6p_{RF}(x)+0.4a_n(x)\right).
\end{equation}
The score is divided into illustrative low, medium, and high bands. The weights (0.6 and 0.4) are heuristic and were not optimized against analyst outcomes or calibrated incident costs. The score should therefore be used to order records for inspection, not interpreted as a probability of compromise.

\subsection{Evaluation metrics and implementation stack}
The experiment reports accuracy, attack precision, attack recall, attack F1-score, false-positive rate (FPR), confusion matrices, and RF ROC-AUC and PR-AUC. For the attack class, precision is $TP/(TP+FP)$, recall is $TP/(TP+FN)$, and F1 is their harmonic mean. FPR is $FP/(FP+TN)$. ROC-AUC and PR-AUC require continuous scores; they are reported for RF probabilities only. AUC values are not calculated for the binary IF and hybrid outputs in the current notebook.

The implementation uses Python 3.11, Jupyter Notebook, pandas, NumPy, matplotlib, joblib, and scikit-learn. The workflow is CPU-based and designed to run on a standard laptop. The notebook saves model artefacts, experiment metadata, feature-importance output, category recall, and risk scores. The complete notebook and setup instructions are included in the accompanying repository package.

\section{Experimental Results}
\subsection{Overall performance}
Table~\ref{tab:metrics} presents the held-out results. RF achieves the strongest overall F1-score (86.22\%) and accuracy (88.63\%). IF has the highest attack precision (96.09\%) and the lowest FPR (1.26\%), but detects only 44.59\% of attack records. The hybrid OR rule achieves the highest recall (86.81\%), but the increase over RF is only 0.13 percentage points. Its precision falls by 0.50 percentage points, F1 decreases by 0.19 percentage points, and FPR increases by 0.43 percentage points.

\begin{table}[t]
\caption{Held-out test performance. Precision, recall, and F1 refer to the attack class.}
\label{tab:metrics}
\centering
\begin{tabular}{lrrrrr}
\toprule
Model & Accuracy & Precision & Recall & F1 & FPR\\
\midrule
Random Forest & 88.63\% & 85.76\% & 86.68\% & 86.22\% & 10.02\%\\
Isolation Forest & 76.51\% & 96.09\% & 44.59\% & 60.91\% & 1.26\%\\
Hybrid OR & 88.43\% & 85.26\% & 86.81\% & 86.03\% & 10.45\%\\
\bottomrule
\end{tabular}
\end{table}

The confusion matrices provide further context. RF produces 50,958 true negatives, 5,675 false positives, 5,251 false negatives, and 34,182 true positives. IF produces 55,918 true negatives, 715 false positives, 21,851 false negatives, and 17,582 true positives. Hybrid OR produces 50,717 true negatives, 5,916 false positives, 5,203 false negatives, and 34,230 true positives. The IF model is conservative, while OR fusion adds 241 true positives compared with RF but also adds 241 false positives. In this test set, the marginal recall improvement therefore comes with a comparable increase in false alarms.

RF obtains ROC-AUC 0.9559 and PR-AUC 0.9479. These values indicate useful ranking ability on this sampled test distribution, but they do not remove the class-prevalence caveat. The AUC values are not directly comparable to the binary predictions of IF and hybrid OR.

\subsection{Feature-level interpretation}
Permutation importance on a sample of up to 3,000 held-out records ranks HTTPS, Max, Number, Time\_To\_Live, and Header\_Length as the five strongest features by mean reduction in attack F1 when shuffled. Other influential features include IAT, Tot sum, Rate, psh\_flag\_number, ack\_count, syn\_count, and Variance. This provides a global view of which flow statistics the RF relies on. It does not establish that any one feature is causally responsible for an attack, nor does it explain an individual alert. Correlation among network features may also affect the ranking.

\subsection{Attack-category results}
The binary model is also assessed by attack category in the held-out test partition. Several flooding and Mirai categories in the test sample receive 100\% recall for both RF and hybrid OR. However, recall is lower for categories such as DictionaryBruteForce (65.86\% RF; 66.05\% hybrid), BrowserHijacking (72.11\%; 72.54\%), and Uploading\_Attack (83.79\%; 84.11\%). The hybrid model's gains are small in these examples. These results should be interpreted carefully because some rare attack categories have only one source CSV in the full data. A category represented by one group cannot be simultaneously represented in both fit and test under strict group separation. The test table therefore reflects only categories and source groups assigned to the test partition; it is not a complete per-category generalization study.

\subsection{Risk prioritization}
The risk score ranks flows using RF attack probability and normalized IF anomaly magnitude. The highest-ranked examples in the executed run have scores around 95--98 and are mostly labelled as attacks in the test set. This demonstrates that the score can create an ordered queue for analyst inspection. However, no human analyst study, response-time measurement, cost-sensitive threshold tuning, or calibrated risk validation was conducted. The score is therefore a prototype triage mechanism rather than a validated operational risk model.

\section{Discussion: Positioning Against Existing Tools}
The framework is a complementary analytics layer, not a replacement for firewalls or mature NIDS/NDR products. Unlike Snort, Suricata, or commercial NDR platforms, it operates on pre-extracted CSV flows and does not capture packets, reassemble streams, match signatures, enrich alerts with threat intelligence, or block traffic. Its narrower potential value is to expose a benign-trained anomaly signal alongside RF probability and rank records for analyst review; it could be integrated with an existing NIDS through a SIEM.

The results caution that hybridization does not automatically improve detection. OR fusion raises recall by only 0.13 percentage points while increasing FPR and reducing F1. This may be acceptable where missed attacks are costly, but undesirable where analyst capacity is limited. Alternative score fusion and context-aware thresholds require empirical testing.

\section{Limitations and Threats to Validity}
The category-capped sample changes class prevalence, so results cannot be extrapolated directly to production. Source-CSV grouping reduces within-file leakage but not dependence across captures or devices; group-based splitting also yields unequal row proportions. Rare categories with one CSV cannot appear in both fit and test. The binary target does not identify attack families. OR fusion and risk weights are heuristics; repeated-seed confidence intervals, significance tests, and cost-sensitive tuning are absent. Permutation importance is global and may be affected by correlated features; rerun it with one worker because the notebook emitted a joblib warning. Live throughput, latency, packet loss, and sensor integration were not measured.

\section{Conclusion and Future Work}
This paper presented an offline hybrid IoT IDS combining RF and benign-trained IF. RF achieved 86.22\% attack F1, compared with 86.03\% for hybrid OR. OR slightly increased recall but also FPR; IF was precise but missed many attacks. Simple OR fusion was not a better overall classifier in this experiment.

Future work includes repeated group splits, device-level holdouts, realistic prevalence, confidence intervals, cost-sensitive fusion, local SHAP explanations, drift monitoring, cross-dataset validation, and live NIDS integration. This remains a research prototype.

\section*{Code Availability}
The complete experiment notebook, environment requirements, and reproduction instructions are provided in the accompanying project repository. The final public repository URL should be inserted here after the author creates and publishes the GitHub repository: \url{https://github.com/[username]/[repository]}. The CICIoT2023 dataset is available from the Canadian Institute for Cybersecurity [1]; the dataset is not redistributed with the code.

\begin{thebibliography}{8}
\bibitem{neto2023} Neto, E.C.P., Dadkhah, S., Ferreira, R., Zohourian, A., Lu, R., Ghorbani, A.A.: CICIoT2023: A real-time dataset and benchmark for large-scale attacks in IoT environment. Sensors 23(13), 5941 (2023). \doi{10.3390/s23135941}
\bibitem{liu2008} Liu, F.T., Ting, K.M., Zhou, Z.-H.: Isolation Forest. In: Proceedings of the 8th IEEE International Conference on Data Mining, pp. 413--422 (2008). \doi{10.1109/ICDM.2008.17}
\bibitem{breiman2001} Breiman, L.: Random forests. Machine Learning 45, 5--32 (2001). \doi{10.1023/A:1010933404324}
\bibitem{fisher2019} Fisher, A., Rudin, C., Dominici, F.: All models are wrong, but many are useful: Learning a variable's importance by studying an entire class of prediction models simultaneously. Journal of Machine Learning Research 20(177), 1--81 (2019).
\end{thebibliography}
\end{document}
